\subsection*{Integracja z aplikacją Meshtastic}
\label{sec:integracja-meshtastic}

Aplikacja MeshTracker nie komunikuje się bezpośrednio z urządzeniami radiowymi - korzysta
z aplikacji Meshtastic zainstalowanej na tym samym telefonie, która pełni rolę pośrednika. 
Wybrano takie podejście, aby nie musieć implementować protokołu Meshtastic i skupić się na 
warstwie prezentacji i funkcjonalności geofencingu. 

\paragraph{Mechanizmy IPC.}

Meshtastic udpostępnia dwa mechanizmy IPC (ang.\ \textit{Inter-Process Communication}):

\begin{itemize}
    \item \textbf{Interfejs AIDL} (ang.\ \textit{Android Interface Definition Language}) -
          umożliwia synchroniczne wywołania metod na działającym serwisie, takie jak pobranie
          listy węzłów czy zapytanie o stan połączenia. MeshTracker korzysta z niego na żądanie,
          głównie przy inicjalizacji i cyklicznym odświeżaniu. \footnote{Na Androidzie jeden proces 
          nie może normalnie uzyskać dostępu do pamięci innego procesu. Aby się komunikować, muszą 
          rozkładać obiekty na elementy, które system operacyjny może zrozumieć, i przekazywać je 
          przez tę granicę. Kod potrzebny do tego przekształcania jest trudny do napisania, dlatego 
          Android wykonuje to za Ciebie za pomocą AIDL.[\url{https://developer.android.com/develop/background-work/services/aidl?hl=pl}]}

    \item \textbf{Broadcasty Androida} - Meshtastic rozsyła komunikaty do wszystkich
          zainteresowanych aplikacji przy każdej zmianie stanu sieci. Pozwala to na reaktywne
          powiadamianie o nowych danych bez aktywnego odpytywania serwisu.  
        \footnote{Aplikacje Androida mogą wysyłać i odbierać komunikaty systemowe oraz pochodzące z innych aplikacji, zgodnie ze wzorem projektowym \textit{publish-subscribe}. \url{https://developer.android.com/guide/components/broadcasts?hl=pl}}
\end{itemize}

MeshTracker korzysta z obu mechanizmów jednocześnie: AIDL do inicjalizacji stanu i cyklicznego
odświeżania, broadcasty do aktualizacji w czasie rzeczywistym. Gwarantuje to
spójność danych niezależnie od tego, czy zdarzenie nastąpi między cyklami odświeżania, czy
podczas startu aplikacji.

Takie podejście do integracji między aplikacjami wprowadza pewną wrażliwość na zmiany API Meshtastic między wersjami.
W szczególności metoda zwracająca ID węzła lokalnego zmieniała swoją nazwę w kolejnych
wydaniach, co wymagało zaimplementowania łańcucha kandydatów:
\texttt{getMyId()} $\to$ \texttt{getMyNodeId()} $\to$ \texttt{myId()} $\to$ \texttt{getOwnerId()}
$\to$ przez \texttt{\seqsplit{getMyNodeInfo().myNodeNum}}. Wersja Meshtastic na telefonie użytkownika
jest poza kontrolą dewelopera MeshTrackera, dlatego integracja musi być możliwie elastyczna i odporna na zmiany w API Meshtastic.

\subsection*{Schemat nawiązywania połączenia.}

Rysunek~\ref{fig:connection-flow} przedstawia sekwencję zdarzeń od uruchomienia
\texttt{MapViewModel} do pojawienia się węzła na mapie.

\begin{figure}[!h]
    \label{fig:connection-flow}
    \centering
    \includegraphics[width=\linewidth]{schematPolaczenia.png}
    \caption{Sekwencja nawiązywania połączenia i przepływ danych}
\end{figure}

\paragraph{Obsługa rozłączenia i błędów.}

Wyróżnione są dwa odrębne scenariusze rozłączenia, które wymagają różnej reakcji aplikacji:

\textbf{Rozłączenie celowe} następuje, gdy użytkownik zamyka aplikację. W \texttt{\seqsplit{MeshServiceManager}} flaga
\texttt{\seqsplit{isIntentionalDisconnect}} jest ustawiana przed wywołaniem \texttt{unbindService()}.
Gdy wywołana jest metoda \texttt{onServiceDisconnected()}, wartość powyższej flagi pozwala
pominąć propagację zdarzenia do warstwy UI, przez co użytkownik nie powinien widzieć komunikatu
o błędzie, skoro rozłączenie było planowane.

\textbf{Rozłączenie nieoczekiwane} - sygnalizuje, że serwis Meshtastic zakończył
działanie bez polecenia MeshTrackera (np.\ na skutek wymuszonego zamknięcia przez system).
W takim przypadku zdarzenie jest propagowane do \texttt{MapViewModel}, który aktualizuje
pasek stanu aplikacji (\texttt{\seqsplit{ConnectionStatusBar}}) odpowiednim komunikatem.

\paragraph{Automatyczna próba ponownego połączenia.}

Odebranie broadcastu \texttt{\seqsplit{ACTION\_MESH\_DISCONNECTED}} \footnote{Broadcasty nie są dostępne w dokumentacji API, jednak można je znaleźć dobrze udokumentowane w repozytorium \url{https://github.com/meshtastic/Meshtastic-Android/blob/main/core/api/src/main/kotlin/org/meshtastic/core/api/MeshtasticIntent.kt}} przy wciąż aktywnym bindingu
z serwisem (tj.\ gdy \texttt{isBound~=~true}) uruchamia w \texttt{\seqsplit{MapViewModel}} procedurę
automatycznego wznawiania połączenia z węzłem lokalnym (\texttt{startReconnecting()}). Działa ona
w pętli, która wykonywana jest aż do skutku lub utraty samego bindingu.

Każda iteracja pętli rozpoczyna się odliczaniem \texttt{\seqsplit{RECONNECT\_INTERVAL\_SECONDS}}~=~30
sekund. W tym czasie stan połączenia ustawiany jest na \texttt{\seqsplit{ConnectionState.Reconnecting}},
który przechowuje liczbę sekund pozostałych do następnej próby. Wartość ta jest co sekundę
odświeżana i wyświetlana użytkownikowi w pasku stanu \texttt{\seqsplit{ConnectionStatusBar}}, dzięki czemu
wie on, kiedy spodziewać się kolejnej próby.

Po upływie odliczania \texttt{MapViewModel} odpytuje serwis przez AIDL o aktualny stan połączenia z węzłem lokalnym.
Jeśli zwrócona wartość to \texttt{STATE\_CONNECTED}, procedura kończy działanie, a aplikacja
przechodzi w stan \texttt{Connected} i wznawia cykliczne odświeżanie węzłów. W przeciwnym razie
licznik prób jest inkrementowany i odliczanie zaczyna się od nowa. Jeśli podczas oczekiwania
binding z serwisem zostanie utracony (\texttt{isConnected()~=~false}), pętla jest przerywana,
a stan przechodzi w \texttt{Disconnected}.

Niezależnie od automatycznych prób, użytkownik może w dowolnej chwili nacisnąć przycisk
\textit{Ponów} widoczny w pasku stanu (rysunek~\ref{fig:connection-status}). Wywołuje to \texttt{retryConnect()},
które anuluje bieżące zadania, resetuje binding i inicjuje nowe połączenie z serwisem
od zera. Jest to przydatne np. gdy użytkownik ręcznie uruchomił aplikację Meshtastic i chce wywołać
natychmiastową próbę połączenia bez czekania na automatyczne odświeżenie.

\begin{figure}[H]
\centering
\begin{tikzpicture}[
    node distance=0.7cm,
    every node/.style={font=\small},
    io/.style={ellipse, draw=black, fill=gray!18,
        text width=4.2cm, minimum height=0.72cm, align=center},
    proc/.style={rectangle, draw=black, fill=blue!7,
        text width=4.2cm, minimum height=0.72cm, align=center},
    dec/.style={diamond, draw=black, fill=orange!15,
        minimum width=3.2cm, minimum height=1.0cm,
        align=center, inner sep=1pt},
    endok/.style={ellipse, draw=black, fill=green!15,
        text width=3.6cm, minimum height=0.72cm, align=center},
    enderr/.style={ellipse, draw=black, fill=red!10,
        text width=3.2cm, minimum height=0.72cm, align=center},
    side/.style={rectangle, draw=black!50, dashed, fill=yellow!10,
        text width=3.5cm, minimum height=0.72cm, align=center},
    arr/.style={->, >=Stealth, thick},
    darr/.style={->, >=Stealth, thick, dashed},
    lbl/.style={font=\scriptsize},
]
% Główny przepływ
\node[io] (recv)  {\small Odebranie \texttt{\seqsplit{ACTION\_MESH\_DISCONNECTED}}\\[1pt]
                   {\scriptsize\texttt{isBound = true}}};
\node[proc,  below=of recv]       (startR) {\texttt{startReconnecting()}};
\node[proc,  below=of startR]     (setR)   {Stan $\to$ \texttt{Reconnecting}};
\node[proc,  below=of setR]       (waitN)  {Odczekaj 30\,s\\[1pt]
                                             {\scriptsize aktualizuj UI co 1\,s}};
\node[dec,   below=0.9cm of waitN] (decB)  {Binding\\utracony?};
\node[proc,  below=0.9cm of decB]  (aidlQ) {Zapytaj serwis AIDL\\[1pt]
                                             o stan połączenia};
\node[dec,   below=0.9cm of aidlQ] (decC)  {\texttt{STATE\_}\\[-1pt]\texttt{CONNECTED?}};
\node[proc,  below=0.9cm of decC]  (incr)  {Inkrementuj\\licznik prób};

% Stany końcowe (prawa kolumna)
\node[enderr, right=2.6cm of decB] (disc) {Stan $\to$ \texttt{Disconnected}};
\node[endok,  right=2.6cm of decC] (conn) {Stan $\to$ \texttt{Connected}\\[1pt]
                                            {\scriptsize wznów odświeżanie}};

% Ścieżka użytkownika (prawa kolumna, u góry)
\node[side, right=2.6cm of startR]  (ponow)   {Użytkownik:\\przycisk ,,Ponów''};
\node[side, below=0.6cm of ponow]   (retryFn) {\texttt{retryConnect()}\\[1pt]
                                               {\scriptsize anuluj, reset bindingu}};

% Strzałki głównego przepływu
\draw[arr] (recv)  -- (startR);
\draw[arr] (startR)-- (setR);
\draw[arr] (setR)  -- (waitN);
\draw[arr] (waitN) -- (decB);
\draw[arr] (decB)  -- node[lbl, above] {Tak} (disc);
\draw[arr] (decB)  -- node[lbl, left=2pt] {Nie} (aidlQ);
\draw[arr] (aidlQ) -- (decC);
\draw[arr] (decC)  -- node[lbl, above] {Tak} (conn);
\draw[arr] (decC)  -- node[lbl, left=2pt] {Nie} (incr);

% Pętla powrotna (lewa strona)
\draw[arr] (incr.west) -- ++(-1.8,0) |- (setR.west);

% Ścieżka przycisku "Ponów"
\draw[darr] (ponow)   -- (retryFn);
\draw[darr] (retryFn) to[bend right=45] (startR.east);
\end{tikzpicture}
\caption{Procedura automatycznego ponownego połączenia}
\label{fig:reconnect-flow}
\end{figure}

\begin{figure}[!h]
    \label{fig:connection-status}
    \centering
    
\includegraphics[width=0.45\linewidth]{nieMoznaPolaczycPasek.jpg}
    \caption{Pasek stanu połączenia - \textit{Ponów}}
\end{figure}

\paragraph{Model danych węzła.}

Obiekt \texttt{NodeInfo} udostępniany przez Meshtastic jest mapowany na wewnętrzną strukturę
danych złożoną z trzech klas: \texttt{MeshNodeInfo}, \texttt{MeshUserInfo} oraz
\texttt{MeshPosition}. 

Klasa \texttt{MeshNodeInfo} przechowuje dane identyfikacyjne i~telemetryczne węzła
(tabela~\ref{tab:meshnodeinfo}).

\begin{table}[!h]
    \caption{Pola klasy \texttt{MeshNodeInfo}}
    \label{tab:meshnodeinfo}
    \centering
    \begin{tabular}{|l|l|p{6.5cm}|}
        \hline
        \textbf{Pole} & \textbf{Typ} & \textbf{Opis} \\
        \hline
        \texttt{num}          & \texttt{Int}           & Unikalny numer węzła w sieci Meshtastic \\
        \texttt{user}         & \texttt{MeshUserInfo} & Dane użytkownika węzła (nazwa, rola, sprzęt) \\
        \texttt{position}     & \texttt{MeshPosition} & Ostatnia znana pozycja GPS \\
        \texttt{snr}          & \texttt{Float}         & SNR ostatnio odebranego pakietu [dB] \\
        \texttt{rssi}         & \texttt{Int}           & RSSI ostatnio odebranego pakietu [dBm] \\
        \texttt{lastHeard}    & \texttt{Int}           & Czas ostatniego odbioru (Unix timestamp [s]) \\
        \texttt{batteryLevel} & \texttt{Int}           & Poziom naładowania baterii [\%] \\
        \texttt{channel}      & \texttt{Int}           & Indeks kanału, na którym węzeł nadaje \\
        \texttt{hopsAway}     & \texttt{Int}           & Liczba przeskoków do węzła w sieci mesh \\
        \hline
    \end{tabular}
\end{table}

Klasa \texttt{MeshUserInfo} zawiera dane konfiguracyjne przypisane do węzła przez jego
właściciela (tabela~\ref{tab:meshuserinfo}).

\begin{table}[!h]
    \caption{Pola klasy \texttt{MeshUserInfo}}
    \label{tab:meshuserinfo}
    \centering
    \begin{tabular}{|l|l|p{6.5cm}|}
        \hline
        \textbf{Pole} & \textbf{Typ} & \textbf{Opis} \\
        \hline
        \texttt{id}            & \texttt{String}  & Identyfikator w formacie \texttt{!XXXXXXXX} \\
        \texttt{longName}      & \texttt{String}  & Pełna nazwa węzła nadana przez użytkownika \\
        \texttt{shortName}     & \texttt{String}  & Skrócona nazwa (do 4 znaków), widoczna na wyświetlaczu węzła \\
        \texttt{hwModelString} & \texttt{String} & Model sprzętu (np.\ \texttt{T\_BEAM}, \texttt{T\_ECHO}, \texttt{T1000\_E}) \\
        \texttt{isLicensed}    & \texttt{Boolean} & Czy węzeł jest zarejestrowany jako amatorska stacja radiowa \\
        \texttt{role}          & \texttt{Int}     & Rola węzła w sieci (CLIENT, ROUTER, TRACKER itp.) \\
        \hline
    \end{tabular}
\end{table}

Klasa \texttt{MeshPosition} przechowuje dane lokalizacyjne węzła
(tabela~\ref{tab:meshposition}). Warto zwrócić uwagę na dwie nieoczywiste konwencje
kodowania stosowane przez Meshtastic: współrzędne geograficzne przechowywane są jako liczby
całkowite (w formacie stopni pomnożonych przez $10^7$) 
\footnote{\url{https://github.com/meshtastic/Meshtastic-Android/blob/19ea164f93bd48dd556d5232d66869b5cb0b2f40/core/model/src/commonMain/kotlin/org/meshtastic/core/model/Node.kt\#L110}} ,
co pozwala na precyzję nawet do około 1 metra bez użycia typów zmiennoprzecinkowych. Natomiast kierunek ruchu (\texttt{groundTrack}) wyrażony jest, z podobnej przyczyny, w setnych stopniach (pomnożonych przez $10^5$).

\begin{table}[!h]
    \caption{Pola klasy \texttt{MeshPosition}}
    \label{tab:meshposition}
    \centering
    \begin{tabular}{|l|l|p{6.5cm}|}
        \hline
        \textbf{Pole} & \textbf{Typ} & \textbf{Opis} \\
        \hline
        \texttt{latitude}        & \texttt{Double} & Szerokość geograficzna [${}^\circ$], po konwersji z $10^7$ \\
        \texttt{longitude}       & \texttt{Double} & Długość geograficzna [${}^\circ$], po konwersji z $10^7$ \\
        \texttt{altitude}        & \texttt{Int}    & Wysokość [m n.p.m.] \\
        \texttt{time}            & \texttt{Int}    & Unix timestamp [s] \\
        \texttt{satellitesInView}& \texttt{Int}    & Liczba widocznych satelitów \\
        \texttt{groundSpeed}     & \texttt{Int}    & Prędkość [m/s] \\
        \texttt{groundTrack}     & \texttt{Int}    & Kierunek ruchu [${}^\circ$](0\,=\,N) \\
        \texttt{precisionBits}   & \texttt{Int}    & Precyzja pozycji (liczba bitów, max 32\,$\approx$\,1\,m) \\
        \hline
    \end{tabular}
\end{table}

Dostępność poszczególnych pól pozycji zależy od konfiguracji węzła nadającego, a konkretnie
od wartości \textit{Position Flags}, czyli bitowego pola opcji włączanych w ustawieniach
radia Meshtastic.\footnote{\url{https://meshtastic.org/docs/configuration/radio/position/\#position-config-values}}
Tabela~\ref{tab:position-flags} zestawia dostępne flagi.

\begin{table}[!h]
    \caption{Flagi \textit{Position Flags} - pola opcjonalne pakietu pozycji}
    \label{tab:position-flags}
    \centering
    \begin{tabular}{|l|p{9cm}|}
        \hline
        \textbf{Flaga} & \textbf{Opis} \\
        \hline
        \texttt{ALTITUDE}           & Dołącza wartość wysokości (jeśli dostępna) \\
        \texttt{ALTITUDE\_MSL}      & Wysokość wyrażona jako MSL (nad poziomem morza) \\
        \texttt{GEOIDAL\_SEPARATION}& Dołącza separację geoidy\footnotemark \\
        \texttt{DOP}                & Dołącza wskaźnik DOP (domyślnie PDOP)\footnotemark{} \\
        \texttt{HVDOP}              & Zamiast PDOP wysyła osobne wartości HDOP i VDOP\footnotemark \\
        \texttt{SATINVIEW}          & Dołącza liczbę widocznych satelitów \\
        \texttt{SEQ\_NO}            & Dołącza numer sekwencyjny pakietu (rośnie co pakiet) \\
        \texttt{TIMESTAMP}          & Dołącza znacznik czasu z rozwiązania GPS \\
        \texttt{HEADING}            & Dołącza kierunek ruchu z rozwiązania GPS \\
        \texttt{SPEED}              & Dołącza prędkość naziemną z rozwiązania GPS \\
        \hline
    \end{tabular}
\end{table}
\addtocounter{footnote}{-2}
\footnotetext{Separacja geoidy (ang.\ \textit{geoidal separation}) to różnica wysokości między
elipsoidą odniesienia (np.\ WGS-84) a geoidą -- powierzchnią ekwipotencjalną odpowiadającą
średniemu poziomowi morza. Niezbędna do przeliczenia wysokości elipsoidalnej GPS na wysokość
nad poziomem morza (MSL).}
\addtocounter{footnote}{1}
\footnotetext{DOP (ang.\ \textit{Dilution of Precision}) to bezwymiarowy wskaźnik jakości geometrii
konstelacji satelitów widocznych w danej chwili. PDOP (ang.\ \textit{Position DOP}) opisuje łączną
precyzję wyznaczania pozycji 3D. Im mniejsza wartość, tym lepsza geometria konstelacji satelitów.}
\addtocounter{footnote}{1}
\footnotetext{HDOP (ang.\ \textit{Horizontal Dilution of Precision}) -- wskaźnik precyzji wyznaczania
pozycji w płaszczyźnie poziomej; VDOP (ang.\ \textit{Vertical Dilution of Precision}) -- analogiczny
wskaźnik dla składowej pionowej.}

W konfiguracji przyjętej w projekcie włączone są flagi \texttt{\seqsplit{ALTITUDE}}, \texttt{\seqsplit{SATINVIEW}},
\texttt{\seqsplit{TIMESTAMP}}, \texttt{\seqsplit{HEADING}} oraz \texttt{\seqsplit{SPEED}}. Flagi są ustawiane per węzeł nadający.
Precyzja ustalania pozycji ustalana jest na poziomie kanału. Ustawiono  maksymalną możliwą 
wartość - 32-bity na kanale \texttt{Reja} Jeśli dane pole nie jest aktywowane, w odebranym 
pakiecie przyjmuje wartość domyślną (zwykle 0), co MeshTracker traktuje jako brak danych. 


